% Generated from locale/tr/content/first-order-logic/completeness/compactness.tex; do not hand-edit.


\iftag{FOL}
      {\olfileid[tr]{fol}{com}{com}}
      {\olfileid[tr]{pl}{com}{com}}

\olsection{Tıkızlık Teoremi}

Tamlık teoreminin önemli sonuçlarından biri tıkızlık teoremidir.
Tıkızlık teoremi, bir !!{sentence}s kümesinin her \emph{sonlu} alt
kümesi sağlanabilir ise bütün kümenin---kendisi sonsuz olsa
bile---sağlanabilir olduğunu söyler. Bu hiç de apaçık değildir. İlk
bakışta, çelişkinin deyim yerindeyse yalnızca sonsuz sayıdaki tümceden
doğduğu çelişkili sonsuz !!{sentence}s kümelerinin var olmasını
engelleyen hiçbir şey görünmez. Tıkızlık teoremi bu olasılığın dışarıda
bırakılabileceğini söyler: Her sonlu alt kümesi sağlanabilir olan
sağlanamaz hiçbir sonsuz !!{sentence}s kümesi yoktur. Tamlık teoremi
gibi bunun da mantıksal gerektirmeyle ilgili bir biçimi vardır: sonsuz
bir !!{sentence}s kümesi bir şeyi mantıksal olarak gerektiriyorsa bunu
zaten sonlu bir alt kümesi gerektirir.

\begin{defn}
  Bir !!{formula}s kümesi~$\Gamma$, her sonlu
  $\Gamma_0\subseteq\Gamma$ sağlanabilir ise ve ancak bu durumda
  \emph{sonlu olarak sağlanabilir}dir.
\end{defn}

\begin{thm}[Tıkızlık Teoremi]
\ollabel{thm:compactness}
Her tümceler kümesi~$\Gamma$ ve her~$!A$ tümcesi için aşağıdakiler
geçerlidir:
\begin{enumerate}
  \item $\Gamma\Entails !A$ ancak ve ancak $\Gamma_0\Entails !A$
    olacak biçimde sonlu bir $\Gamma_0\subseteq\Gamma$ varsa.
  \item $\Gamma$ sağlanabilirdir ancak ve ancak sonlu olarak
    sağlanabilir ise.
\end{enumerate}
\end{thm}

\begin{proof}
(2)'yi ispatlayalım. $\Gamma$ sağlanabilirse, her $!A\in\Gamma$ için
$\iftag{FOL}{\Sat{M}}{\pSat{v}}{!A}$ olacak biçimde
  \iftag{FOL}{!!a{structure}~$\Struct{M}$}{!!a{valuation}~$\pAssign{v}$}
  vardır. Bu $\iftag{FOL}{\Struct{M}}{\pAssign{v}}$ elbette~$\Gamma$'nın
her sonlu alt kümesini de sağlar; dolayısıyla $\Gamma$ sonlu olarak
sağlanabilirdir.

Şimdi $\Gamma$'nın sonlu olarak sağlanabilir olduğunu varsayalım.
O zaman her sonlu $\Gamma_0\subseteq\Gamma$ sağlanabilirdir. Sağlamlık
uyarınca
(\iftag{FOL}{%
  \tagrefs{prfAX/{fol:axd:sou:cor:consistency-soundness},
    prfSC/{fol:seq:sou:cor:consistency-soundness},
    prfND/{fol:ntd:sou:cor:consistency-soundness},
    prfTab/{fol:tab:sou:cor:consistency-soundness}}}{%
  \tagrefs{prfAX/{pl:axd:sou:cor:consistency-soundness},
    prfSC/{pl:seq:sou:cor:consistency-soundness},
    prfND/{pl:ntd:sou:cor:consistency-soundness},
    prfTab/{pl:tab:sou:cor:consistency-soundness}}})
her sonlu alt küme tutarlıdır. Ardından
\iftag{FOL}{%
  \tagrefs{prfAX/{fol:axd:ptn:prop:proves-compact},
    prfSC/{fol:seq:ptn:prop:proves-compact},
    prfND/{fol:ntd:ptn:prop:proves-compact},
    prfTab/{fol:tab:ptn:prop:proves-compact}}}{%
  \tagrefs{prfAX/{pl:axd:ptn:prop:proves-compact},
    prfSC/{pl:seq:ptn:prop:proves-compact},
    prfND/{pl:ntd:ptn:prop:proves-compact},
    prfTab/{pl:tab:ptn:prop:proves-compact}}}
uyarınca $\Gamma$'nın kendisi de tutarlı olmalıdır. Tamlık
(\olref[cth]{thm:completeness}) uyarınca $\Gamma$ tutarlı olduğundan
sağlanabilirdir.
\end{proof}

\tagprob{FOL}
\begin{prob}
\olref[fol][com][com]{thm:compactness} teoreminin (1). maddesini
ispatlayın.
\end{prob}
\tagendprob

\tagprob{notFOL}
\begin{prob}
\olref[pl][com][com]{thm:compactness} teoreminin (1). maddesini
ispatlayın.
\end{prob}
\tagendprob

\iftag{FOL}{%
\begin{ex}
Bir~$\Gamma$ kuramının her~$\Struct{M}$ modelinde her~$t$ terimi elbette
$\Domain{M}$ içindeki !!a{element} öğesini gösterir. Tersinin de doğru
olduğunu, yani $\Domain{M}$ içindeki her !!{element} öğesinin bir terim
tarafından gösterildiğini güvence altına alabilir miyiz? Başka bir
deyişle, bütün modelleri örtülmüş olan~$\Gamma$ kuramları var mıdır?
Tıkızlık teoremi, $\Gamma$'nın sonsuz modelleri varsa bunun mümkün
olmadığını gösterir. Şöyle görelim: $\Struct{M}$,~$\Gamma$'nın sonsuz
bir modeli ve $c$,~$\Gamma$ dilinde bulunmayan !!a{constant} olsun.
$\Delta$,~$\Gamma$'nın $\Lang L$ dilindeki her~$t$ terimi için bütün
$\eq/[c][t]$ tümcelerinin kümesi olsun; yani
  \[
  \Delta = \Setabs{\eq/[c][t]}{t \in \Trm[L]}.
  \]
$\Gamma\cup\Delta$ kümesinin sonlu bir alt kümesi, $\Gamma'\subseteq
\Gamma$ ve $\Delta'\subseteq\Delta$ olmak üzere
$\Gamma'\cup\Delta'$ biçiminde yazılabilir. $\Delta'$ sonlu olduğundan
yalnızca sonlu sayıda terim içerebilir. $a\in\Domain{M}$ olsun ve $a$,
$\Domain{M}$ içinde bunların hiçbirinin göstermediği !!a{element} olsun;
!!{structure}~$\Struct{M'}$ yapısını, $\Struct{M}$ ile aynı fakat ayrıca
$\Assign{c}{M'}=a$ olacak
biçimde tanımlayalım. $\Delta'$ içinde geçen bütün~$t$ terimleri için
$a\neq\Value{t}{M}$ olduğundan $\Sat{M'}{\Delta'}$. Ayrıca
$\Sat{M}{\Gamma}$, $\Gamma'\subseteq\Gamma$ ve $c$,~$\Gamma$ içinde
geçmediğinden $\Sat{M'}{\Gamma'}$. Böylece $\Gamma\cup\Delta$'nın her
sonlu $\Gamma'\cup\Delta'$ alt kümesi için
$\Sat{M'}{\Gamma'\cup\Delta'}$ olur. Demek ki $\Gamma\cup\Delta$'nın
her sonlu alt kümesi sağlanabilirdir. Tıkızlık gereği
$\Gamma\cup\Delta$'nın kendisi de sağlanabilirdir. Dolayısıyla
$\Sat{M}{\Gamma\cup\Delta}$ olan modeller vardır. Böyle her~$\Struct{M}$,
$\Gamma$'nın bir modelidir fakat örtülmüş değildir; çünkü~$\Lang L$
dilinin bütün~$t$ terimleri için
$\Value{c}{M}\neq\Value{t}{M}$.
\end{ex}

\begin{ex}
!!a{language}~$\Lang L$; !!{predicate}~$<$, !!{constant}s~$\Obj{0}$ ve
$\Obj{1}$ ile !!{function}s~$+$, $\times$ ve~$-$ simgelerini içersin.
$\Gamma$, bütün !!{sentence}s ifadelerinin kümesi olsun: bu ifadeler
bu !!{language} içinde olup evreni~$\Rat$ ve yorumları olağan olan
!!{structure}~$\Struct{Q}$ içinde doğrudur. $\Gamma$, rasyonel sayılar
hakkında doğru olan bütün~$\Lang L$ !!{sentence}s ifadelerinin kümesidir.
Elbette~$\Rat$ içinde (hatta~$\Real$ içinde), $0$'dan
büyük fakat her $k\in\PosInt$ için $1/k$'dan küçük hiçbir~$r$ sayısı
yoktur. Böyle bir sayı var olsaydı \emph{sonsuz küçük} olurdu: sıfırdan
farklı ama sonsuz ölçüde küçük. Tıkızlık teoremi,~$\Gamma$'nın sonsuz
küçüklerin var olduğu modellerinin bulunduğunu göstermek için
kullanılabilir. Dilimizde bölme için !!a{function} yoktur (sıfıra bölme
tanımsızdır ve !!{function}s tümel fonksiyonlarla yorumlanmalıdır).
Bununla birlikte $r<1/k$ olduğunu yine de ifade edebiliriz; çünkü bu,
ancak ve ancak $r\cdot k<1$ ise geçerlidir. $c$ yeni bir !!{constant}
olsun ve $\Delta$ şu küme olsun:
\[
\{\Obj{0}<c\}\cup
\Setabs{c\times\num{k}<\Obj{1}}{k\in\PosInt}
\]
(burada $\num{k}=(\Obj{1}+(\Obj{1}+\dots+(\Obj{1}+\Obj{1})\dots))$
ve ifadede $k$ tane~$\Obj{1}$ vardır). $\Delta$'nın herhangi bir sonlu
$\Delta_0$ alt kümesi için öyle bir~$K$ vardır ki $\Delta_0$ içindeki
bütün $c\times\num{k}<\Obj{1}$ !!{sentence}s ifadelerinde $k<K$ olur.
$\Struct{Q}$ yapısını $\Assign{c}{Q'}=1/K$ olacak biçimde
$\Struct{Q'}$ yapısına genişletirsek her sonlu $\Gamma_0\subseteq
\Gamma$ için $\Sat{Q'}{\Gamma_0\cup\Delta_0}$ olur. Dolayısıyla
$\Gamma\cup\Delta$ sonlu olarak sağlanabilirdir (Alıştırma: bunu
ayrıntılı olarak ispatlayın). Tıkızlık gereği $\Gamma\cup\Delta$
sağlanabilirdir. $\Gamma\cup\Delta$'nın her~$\Struct{S}$ modeli,
$\Assign{c}{S}$ sonsuz küçük öğesini içerir.
\end{ex}
}{}

\tagprob{FOL}
\begin{prob}
Aritmetiğin standart modeli~$\Struct{N}$ içinde her $\num{n}<x$
formülünü sağlayan hiçbir $k\in\Domain{N}$ !!{element} öğesi yoktur
(burada $\num{n}$, $n$ tane $\prime$ içeren
$\Obj{0}^{\prime\dots\prime}$ ifadesidir). Tıkızlık teoremini kullanarak,
aritmetik dilinde standart~$\Struct{N}$ modelinde doğru olan bütün
!!{sentence}s ifadelerinin !!a{structure}~$\Struct{N'}$ içinde de doğru
olduğunu; bu yapının her $\num{n}<x$ formülünü \emph{sağlayan}
!!a{element} içerdiğini gösterin.
\end{prob}
\tagendprob

\iftag{FOL}{
\begin{ex}
!!{identity} içeren birinci dereceden mantığın, !!{domain} boyutunun
belirli bir alt sınırı olması gerektiğini ifade edebildiğini biliyoruz:
``en az $n$ farklı nesne vardır'' diyen $!A_{\ge n}$ tümcesi yalnızca
$\Domain{M}$ içinde en az~$n$ nesne bulunan yapılarda doğrudur. Bu yüzden
  \[
  \Delta=\Setabs{!A_{\ge n}}{n\ge1}
  \]
alırsak $\Delta$'nın her modeli sonsuz olmalıdır. Böylece bir kurama
$\Delta$ ekleyerek onun yalnızca sonsuz modellere sahip olmasını
güvence altına alabiliriz: $\Gamma\cup\Delta$'nın modelleri tam olarak
$\Gamma$'nın sonsuz modelleridir.

Demek ki birinci dereceden mantık sonsuzluğu ifade edebilir. Fakat
tıkızlık teoremi sonluluğu ifade edemediğini gösterir. Aksine, bir
$\Lambda$ tümceler kümesinin tam olarak sonlu !!{structure}s içinde
sağlandığını varsayalım. O zaman $\Delta\cup\Lambda$ sonlu olarak
sağlanabilirdir. Neden? $\Delta'\subseteq\Delta$ ve
$\Lambda'\subseteq\Lambda$ olmak üzere sonlu bir
$\Delta'\cup\Lambda'\subseteq\Delta\cup\Lambda$ alalım.
$\Delta'$ boşsa $n=1$ alalım; değilse
$!A_{\ge n}\in\Delta'$ olacak biçimdeki en büyük sayıya~$n$ diyelim.
Bütün sonlu yapılarda sağlanan~$\Lambda$, sonlu fakat $\ge n$
!!{element}s içeren bir~$\Struct{M}$ modeline sahiptir. O zaman
$\Sat{M}{\Delta'\cup\Lambda'}$. Tıkızlık gereği
$\Delta\cup\Lambda$'nın sonsuz bir modeli vardır; bu,~$\Lambda$'nın
yalnızca sonlu !!{structure}s içinde sağlandığı varsayımıyla çelişir.
\end{ex}
}{}

